\documentclass[a4paper,12pt]{report}

\usepackage{alltt, fancyvrb, url}
\usepackage{graphicx}
\usepackage[utf8]{inputenc}
\usepackage{float}
\usepackage{xcolor}
\usepackage{hyperref}

\usepackage[italian]{babel}

\usepackage[italian]{cleveref}

\title{Relazione DungeonSQL}
\author{Andrea Maria Castronovo, \\ Carmelo Falsone}
\date{\today}


\begin{document}

\maketitle

\tableofcontents

\chapter{Analisi}

\vspace{-20pt}

\section{Introduzione}

L'applicazione si pone l'obiettivo di facilitare la gestione e l'archiviazione di ciò 
che riguarda una o più campagne del gioco di ruolo Dungeons and Dragons, occupandosi 
principalmente di gestire più facilmente le parti più "macchinose" del gioco, 
lasciando un po' più spazio alla parte di role play dei giocatori. 
%
Ad esempio fornirà ai giocatori dei modi semplici per tenere traccia dei loro bonus,
quanti slot magici hanno utilizzato o possono ancora utilizzare oppure la vita rimasta 
dopo un combattimento.

Una \textbf{Campagna} (ovvero il mondo di gioco) necessita di un \textit{nome}, una \textit{descrizione}
e la presenza di un gruppo di \textbf{Giocatori} 
(utenti) di cui uno (il creatore) ricoprirà il ruolo di \textbf{Dungeon Master} (DM).
%
Una campagna è composta da \textbf{Sessioni di gioco} che si tengono in determinati \textit{giorni} 
concordati dai giocatori. Il DM potrà appuntare le cose che accadono nelle sessioni in 
un \textbf{Diario} della campagna.
%
Ogni giocatore potrà creare delle campagne assumendo il ruolo di DM e inviterà altri 
giocatori a partecipare con un loro \textbf{Personaggio}.
%
Ogni utente potrà partecipare ad una campagna con un solo personaggio, il quale dovrà 
essere creato oppure selezionato (ed eventualmente riadattato) da quelli preparati 
precedentemente. 
%
Il DM può usare qualsiasi suo personaggio come personaggio non giocante (NPC) 
e può creare dei combattimenti all'interno della campagna selezionando i personaggi coinvolti che vuole 
oppure qualsiasi NPC o \textbf{Mostro} a sua scelta. 
%
Un combattimento è diviso in round e turni,
questi ultimi ordinati in base ad un tiro iniziativa di ogni partecipante. 
La creazione di un personaggio comprende diversi aspetti stabiliti in un manuale scelto.
%
Per i fini di questo progetto sarà tenuta in considerazione una parte del manuale della 
5° edizione di Dungeons and Dragons. Un personaggio sarà creato scegliendo una \textbf{Razza}
(divisa anche in \textbf{Sottorazze}), una \textbf{Classe} 
(più avanti sviluppata in \textbf{Sottoclasse}) e tutto ciò che il manuale dice sulla classe 
(ad esempio i \textbf{Tratti}). 
%
Un personaggio ha dei "punti caratteristica" da dover distribuire nelle 
seguenti caratteristiche principali: 
\textit{forza}, \textit{costituzione}, \textit{destrezza}, \textit{intelligenza}, \textit{saggezza} e 
\textit{carisma}, tenendo conto degli incrementi forniti ad esempio dalla razza; 
%
inoltre ha anche delle \textit{monete}, un inventario composto da \textbf{Oggetti} e un 
equipaggiamento composto da \textbf{Armi} e \textbf{Armature}. 
%
Esistono inoltre vari \textbf{Incantesimi} che possono essere utilizzati
dai personaggi, che però dovranno avere adeguati requisiti per utilizzarli (ad esempio 
un incantesimo può essere appreso solo da determinate classi o necessita un determinato livello di una classe).


\end{document}